\epanalipsi{7}

\tmima{Λεξικό του κεφαλαίου}

\renewcommand{\arraystretch}{1.5}
\begin{tabular*}{\linewidth}{@{\extracolsep{\fill}} >{\raggedright\arraybackslash}p{0.28\linewidth} >{\raggedright\arraybackslash}p{0.62\linewidth}}
\toprule
\textbf{Λέξη ή τύπος} & \textbf{Με μία πρόταση}\\
\midrule
ρητοί & Φυσικοί, κλάσματα και δεκαδικοί μαζί με τους αντίστοιχους αρνητικούς.\\
$|\alpha|$ & Η απόσταση του $\alpha$ από το $0$ στον άξονα: ποτέ αρνητική.\\
αντίθετοι & Ίδια απόλυτη τιμή, ετερόσημοι: $\alpha$ και $-\alpha$, με άθροισμα $0$.\\
σύγκριση & Μεγαλύτερος είναι ο δεξιότερος στον άξονα.\\
$\alpha-\beta$ & $=\alpha+(-\beta)$: προσθέτω τον αντίθετο.\\
γινόμενο, πηλίκο & Ομόσημοι: θετικό. Ετερόσημοι: αρνητικό. Όχι διαίρεση με το $0$.\\
περιοδικός & Δεκαδικός με ψηφία που επαναλαμβάνονται, π.χ. $\frac{1}{3}=0{,}\overline{3}$.\\
$\alpha^{\nu}$, $\alpha^0$, $\alpha^{-\nu}$ & $\alpha\cdots\alpha$, $1$ και $\frac{1}{\alpha^\nu}$.\\
$\alpha\cdot10^{\nu}$ & Τυποποιημένη μορφή: $1\le\alpha<10$, $\nu$ ακέραιος.\\
\bottomrule
\end{tabular*}

\tmima{Ασκήσεις επανάληψης}

\ex Εξέφρασε με ρητό αριθμό: α) $8$ βαθμοί κάτω από το μηδέν \quad β) χρέος $25$ € \quad γ) βάθος $40$ μ. κάτω από την επιφάνεια της θάλασσας.

\ex Ποιοι από τους $-3$, $\frac{4}{2}$, $0{,}5$, $-\frac{7}{7}$, $0$ είναι ακέραιοι;

\ex Βρες: α) $|-7{,}5|$ \quad β) $\left|\frac{3}{4}\right|$ \quad γ) τον αντίθετο του $-\frac{2}{3}$ \quad δ) τους $\alpha$ με $|\alpha|=4$.

\ex Γράψε σε αύξουσα σειρά: $-2{,}5$, $0$, $-\frac{1}{2}$, $3$, $-3$.

\ex Υπολόγισε: α) $(-7)+(+12)$ \quad β) $(-3{,}5)+(-1{,}5)$ \quad γ) $\left(-\frac{1}{2}\right)+\left(+\frac{1}{3}\right)$ \quad δ) $(+8)+(-8)$.

\ex Υπολόγισε: α) $5-(-3)$ \quad β) $(-4)-(+6)$ \quad γ) $(-2{,}5)-(-2{,}5)$ \quad δ) $\left(-\frac{1}{2}\right)-\frac{1}{4}$.

\ex Υπολόγισε το αλγεβρικό άθροισμα $-12+7-(-5)-(3-9)$.

\ex Υπολόγισε: α) $(-6)\cdot(-7)$ \quad β) $(+4)\cdot(-2{,}5)$ \quad γ) $\left(-\frac{1}{2}\right)(-4)(-3)$ \quad δ) $(-2)\cdot(5-8)$.

\ex Υπολόγισε: α) $(-48):(-6)$ \quad β) $15:(-5)$ \quad γ) $\left(-\frac{3}{4}\right):\left(-\frac{9}{2}\right)$ \quad δ) $0:(-7)$.

\ex Γράψε σε δεκαδική μορφή τα $\frac{5}{6}$ και $\frac{7}{20}$ και ως κλάσμα τον $0{,}\overline{18}$.

\ex Υπολόγισε: α) $(-2)^5$ \quad β) $-2^4$ \quad γ) $\left(\frac{1}{3}\right)^{-2}$ \quad δ) $5^0+2^{-3}$.

\ex Υπολόγισε με ιδιότητες: α) $2^5\cdot2^{-3}$ \quad β) $(3^{-1})^2$ \quad γ) $10^4:10^6$.

\ex Γράψε σε τυποποιημένη μορφή τα $0{,}00032$ και $45\,000\,000$, και υπολόγισε $(4\cdot10^{3})(5\cdot10^{-5})$.

\ex \textbf{Η πτώση.} Η θερμοκρασία πέφτει από $+4$° σε $-5$° σε $3$ ώρες. Ποιος είναι ο ρυθμός μεταβολής ανά ώρα;

\ex \textbf{Ερωτήσεις αυτοαξιολόγησης (πράξεις).} Σωστό ή λάθος; α) $|-5|=-5$. \quad β) $(-3)+(-4)=-7$. \quad γ) $5-(-2)=3$. \quad δ) $(-2)(-3)(-4)=-24$. \quad ε) $12:(-4)=3$. \quad στ) Ο αντίθετος του $0$ είναι το $0$.

\ex \textbf{Ερωτήσεις αυτοαξιολόγησης (δυνάμεις, δεκαδικοί).} Σωστό ή λάθος; α) $(-2)^0=1$. \quad β) $2^{-1}=-2$. \quad γ) Το $\frac{1}{7}$ είναι περιοδικός δεκαδικός. \quad δ) $10^{-3}>10^{-4}$. \quad ε) $(\alpha+\beta)^2=\alpha^2+\beta^2$. \quad στ) $3{,}2\cdot10^{-2}=0{,}032$.

\tmima{Ο χάρτης της επανάληψης}

Μετά τις λύσεις, δες σε ποιες ασκήσεις δυσκολεύτηκες. Ο πίνακας σου λέει ποια παράγραφο να ξαναδιαβάσεις.

\renewcommand{\arraystretch}{1.5}
\begin{tabular*}{\linewidth}{@{\extracolsep{\fill}} >{\raggedright\arraybackslash}p{0.38\linewidth} >{\raggedright\arraybackslash}p{0.52\linewidth}}
\toprule
\textbf{Αν δυσκολεύτηκες στις} & \textbf{ξαναδιάβασε τις παραγράφους}\\
\midrule
Ε7.1, Ε7.2 & 7.1 (ρητοί, άξονας)\\
Ε7.3, Ε7.4 & 7.2 (απόλυτη τιμή, αντίθετοι, σύγκριση)\\
Ε7.5 & 7.3 (πρόσθεση)\\
Ε7.6, Ε7.7 & 7.4 (αφαίρεση, παρενθέσεις)\\
Ε7.8 & 7.5 (πολλαπλασιασμός)\\
Ε7.9, Ε7.14 & 7.6 (διαίρεση)\\
Ε7.10 & 7.7 (δεκαδική μορφή)\\
Ε7.11, Ε7.12 & 7.8 και 7.9 (δυνάμεις)\\
Ε7.13 & 7.10 (τυποποιημένη μορφή)\\
Ε7.15, Ε7.16 & 7.1 έως 7.10\\
\bottomrule
\end{tabular*}
